\abstract
RAM-USB è un servizio di backup multi-utente, geo-distribuito e accessibile da remoto, progettato
secondo tre principi di sicurezza: zero-knowledge, nessun componente lato server accede mai al
contenuto dei backup in chiaro; zero-trust, nessun componente si fida implicitamente dei dati ricevuti
da un altro; defense-in-depth, ogni strato rivalida indipendentemente l'input ricevuto. Il progetto nasce 
come caso di studio approfondito sulla progettazione di sistemi di backup distribuiti.

Il sistema è un'architettura a microservizi n-tier: ogni registrazione o login attraversa quattro
processi distinti, ciascuno dei quali rivalida l'input senza fidarsi dello strato precedente, mentre
l'accesso allo spazio di backup, isolato per utente tramite chroot dinamico, resta possibile solo dopo
l'autenticazione, su una rete privata mesh che applica confini di raggiungibilità distinti per ogni fase
del ciclo di vita dell'utente. I requisiti sono raccolti in un unico documento tracciabile fino al
codice, con un modello a spirale per l'analisi e un modello a V per l'implementazione.

L'analisi critica delle proprietà di sicurezza quantifica il costo di un attacco offline alle password,
giorni per la lunghezza minima ammessa, ordini di grandezza superiori per password più lunghe, e
verifica componente per componente che la compromissione isolata di un servizio non si propaghi agli
altri. Ogni scelta architetturale è messa a bilancio con il prezzo che comporta, dalla latenza
introdotta dalla separazione in più processi alla complessità della rivalidazione indipendente.

Il sistema è stato distribuito e verificato end-to-end su un'infrastruttura reale. 
I suoi limiti del sisitema sono dichiarati insieme agli sviluppi che potrebbero mitigarli.
